% Lösungen Einheit 2 von 5 – Einsetzungsverfahren (zusammenbau v0.1)
\einheitenkopf{Einheit 2 von 5 · Einsetzungsverfahren}

%% TODO Lösung der Erklärzeile 16g) fehlt in der Bank
\erg{16}{a) Gleichung II \quad b) $x = 3$, $y = 7$; Lösung $(3|7)$ \quad c) $x = 4$, $y = 7$; Lösung $(4|7)$ \quad d) $x = 4$, $y = 10$; Lösung $(4|10)$ \quad e) $x = 3$, $y = 12$; Lösung $(3|12)$ \quad f) $x = 3$, $y = 11$; Lösung $(3|11)$ \quad g) \ldots \quad h) $x = 4$, $y = 1$; Lösung $(4|1)$ \quad i) $y = 3$, $x = 1$; Lösung $(1|3)$ \quad j) $x = 3$, $y = 5$; Lösung $(3|5)$ \quad k) $x = 5$, $y = 9$; Lösung $(5|9)$ \quad l) $x = 2{,}40$, $y = 1{,}10$ \quad m) $x = -2$, $y = 3$; Lösung $(-2|3)$ \quad n) $x = 3$, $y = 5$; Lösung $(3|5)$ \quad o) $x = 5$, $y = 3$; Probe: $5 + 9 = 14$ (wA), $10 - 3 = 7$ (wA); Lösung $(5|3)$ \quad p) $y = 54 - x$ in II: $7x + 216 - 4x = 316{,}20$, $x = 33{,}40$; $y = 20{,}60$. Ein Helm kostet $33{,}40\,€$, ein Schloss $20{,}60\,€$.}
\erg{17}{a) $x = \frac{5}{3}$, $y = \frac{10}{3}$}
\erg{18}{a) Der Term aus I muss in II eingesetzt werden: $2x + 3(6 - x) = 14$, also $x = 4$, $y = 2$. \quad b) Die $2$ gilt für beide Glieder: $4x + 2x + 6 = 30$, also $x = 4$, $y = 7$. \quad c) Die Lösung ist ein Zahlenpaar; es fehlt $y = 3 \cdot 3 - 1 = 8$. Lösung: $x = 3$, $y = 8$. \quad d) In II wird $y$ durch den Term $2x + 3$ ersetzt, der nur $x$ enthält; also kommt $y$ nicht mehr vor.}
